\clearpage
\VerdierPhysicalPageAnchor{0048}%
\VerdierSemanticAnchor{ch1-1-4-4}
\begingroup
\emergencystretch=1em
\noindent\textbf{定義 1.4.4.}\quad
$(\mathcal C,G,T,c)$ と $(\mathcal C',G,T',c')$ を二つの $G$-圏とし、
$\mathcal C'$ は加法圏であるとする。また
$\alpha\in Z^2(G,\mathbb Z/2\mathbb Z)$ を正規化された非斉次余サイクルとする。
関手

\VerdierSemanticAnchor{ch1-1-4-4-functor}
\[
  F\colon\mathcal C\longrightarrow\mathcal C'
\]
と関手同型

\VerdierSemanticAnchor{ch1-1-4-4-twisting-data}
\[
  p(\sigma)\colon FT(\sigma)\longrightarrow T'(\sigma)F,
  \qquad \sigma\in G
\]
からなる対象 $(F,p)$ を、\emph{余サイクル $\alpha$ で捩られた $G$-関手}
という。ただし各 $\sigma_1,\sigma_2\in G$ に対して図式 (1.4.3.11) が
$\alpha(\sigma_1,\sigma_2)$-可換であることを要求する。
\par
\endgroup

\VerdierSemanticAnchor{ch1-1-4-4-correspondence}
したがって拡大 (1.4.3.4) の正規化された切断 $s$ を選ぶと、$F$ 上に符号群の
標準作用を誘導する $H$-関手
\[
  (F,m)\colon(\mathcal C,H,T_f,c_f)
  \longrightarrow(\mathcal C',H,T'_f,c'_f)
\]
と、切断 $s$ が定める余サイクル $\alpha_s$ で捩られた $G$-関手との間に
全単射が得られる。

\VerdierSemanticAnchor{ch1-1-4-5}
\noindent\textbf{1.4.5.}\quad
(1.4.3) の記法を用いる。

\VerdierSemanticAnchor{ch1-1-4-5-eq1}
\begin{equation}
  \begin{gathered}
    (F_1,m_1),(F_2,m_2)\colon
    (\mathcal C,H,T_f,c_f)\\[-0.2em]
    \longrightarrow(\mathcal C',H,T'_f,c'_f)
  \end{gathered}
  \tag{1.4.5.1}
\end{equation}
を符号群の標準作用を誘導する二つの $H$-関手とし、

\VerdierSemanticAnchor{ch1-1-4-5-eq2}
\begin{equation}
  u\colon(F_1,m_1)\longrightarrow(F_2,m_2)
  \tag{1.4.5.2}
\end{equation}
を $H$ の作用と両立する関手の射とする。

\VerdierSemanticAnchor{ch1-1-4-5-sections}
完全列 (1.4.3.4) の写像 $f$ の正規化された二つの切断を $s_1,s_2$ とする
（$s_1(e)=s_2(e)=e$）。対応する正規化された余サイクルを
$\alpha_{s_1},\alpha_{s_2}\in Z^2(G,\mathbb Z/2\mathbb Z)$ とし、写像
$\beta\colon G\to\mathbb Z/2\mathbb Z$ を

\VerdierSemanticAnchor{ch1-1-4-5-eq3}
\begin{equation}
  \beta(\sigma)=s_2(\sigma)s_1(\sigma)^{-1},
  \qquad \sigma\in G
  \tag{1.4.5.3}
\end{equation}
で定める。

\VerdierSemanticAnchor{ch1-1-4-5-induced}
(1.4.3) により、関手 $(F_1,m_1)$ と $(F_2,m_2)$ は、それぞれ余サイクル
$\alpha_{s_1}$ と $\alpha_{s_2}$ で捩られた $G$-関手 $(F_1,p_1)$ と
$(F_2,p_2)$ を定める。基礎にある関手の間の射 $u\colon F_1\to F_2$ が
次の性質をもつことは直ちに確かめられる：

\VerdierSemanticAnchor{ch1-1-4-5-property}
